\documentclass{article}
\usepackage{siunitx}
\usepackage{amsmath}
\usepackage{geometry}% 1-inch margins
\geometry{letterpaper, margin=1in}% 1-inch margins
\title{Help Document: EMCH 290 HW \#2}
\author{Jacob C. Vaught}
\begin{document}
\maketitle

\section*{Introduction}
This help document aims to provide you with the necessary equations and example problems to assist with your homework questions on topics such as kinetic energy, potential energy, polytropic processes, and thermodynamics. Each section will focus on a specific topic and provide example problems distinct from the ones in your solution manual.

\section*{Kinetic and Potential Energy}

\subsection*{Key Equations}
\begin{align*}
    \text{Kinetic Energy (KE)} &= \frac{1}{2}mv^2 \\
    \text{Potential Energy (PE)} &= mgh \\
    \text{Work-Energy Principle} & : W = \Delta KE \\
    \text{Energy Conservation} &: KE_i + PE_i = KE_f + PE_f
\end{align*}

\subsection*{Solutions}

\subsubsection*{Example 1: Rotating Disc}

A rotating disc has a moment of inertia \(I = \SI{150}{kg.m^2}\) and is spinning at an angular velocity of \(\omega = \SI{10}{rad/s}\). Calculate the kinetic energy.

\begin{align*}
    KE &= \frac{I \omega^2}{2} \\
    &= \frac{150 \times (10)^2}{2} \\
    &= \SI{7500}{J} = \SI{7.5}{kJ}
\end{align*}

\section*{Polytropic Processes}

\subsection*{Key Equations}

\begin{align*}
    \text{Polytropic Equation} &: pv^n = \text{constant} \\
    \text{Work done} &: W = m \left( \frac{p_2v_2 - p_1v_1}{1-n} \right)
\end{align*}

\subsection*{Solutions}

\subsubsection*{Example 2: Gas Compression}

A gas is compressed from a specific volume of \(v_1 = \SI{1}{m^3/kg}\) to \(v_2 = \SI{0.4}{m^3/kg}\) while the pressure changes from \(p_1 = \SI{100}{kPa}\) to \(p_2 = \SI{300}{kPa}\). Calculate the polytropic exponent \(n\).

\begin{align*}
    n &= \frac{\ln\left(\frac{p_1}{p_2}\right)}{\ln\left(\frac{v_2}{v_1}\right)} \\
    &= \frac{\ln\left(\frac{100}{300}\right)}{\ln\left(\frac{0.4}{1}\right)} \\
    &= \SI{1.22}{}
\end{align*}

\section*{Thermodynamics and Heat Pumps}

\subsection*{Key Equations}

\begin{align*}
    \text{Total Energy Transfer} & : Q = Q_{\text{in}} + Q_{\text{out}} \\
    \text{Coefficient of Performance} & : \gamma = \frac{Q_{\text{out}}}{W}
\end{align*}

\subsection*{Solutions}

\subsubsection*{Example 3: Coefficient of Performance}

A heat pump transfers \SI{30}{kJ} of energy from the external environment and utilizes \SI{10}{kJ} of electrical work. Calculate the coefficient of performance.

\begin{align*}
    \gamma &= \frac{Q_{\text{out}}}{W} \\
    &= \frac{30}{10} \\
    &= \SI{3}{}
\end{align*}

\end{document}
